%%%%%%%%%%%%%%%%%%%%%%%%%%%%%%%%%%%%%%%%%%%%%%%%%%%%%%%%%%%%%%%%%%%%%%%%%%%%%%%
% SECTION 4.X: FROM EXPERIMENTAL EVIDENCE TO CALIBRATED PREDICTION
% Why This Framework Differs Fundamentally from LLM-Based Approaches
% EXPANDED VERSION with worked examples, objections, and rhetorical sharpening
%%%%%%%%%%%%%%%%%%%%%%%%%%%%%%%%%%%%%%%%%%%%%%%%%%%%%%%%%%%%%%%%%%%%%%%%%%%%%%%

% =============================================================================
% Chapter 4x: Calibration not Simulation
% =============================================================================
%
% Version: 1.0 (January 2026)
%
% Purpose: LLM Monte Carlo method
%
% Primary Appendix: See appendix references below
% Secondary Appendices: G (Glossary), F (Examples)
%
% Prerequisites: None
% Leads to: Chapter 1
%
% Chapter Type: B (Foundation)
% =============================================================================

%%%%%%%%%%%%%%%%%%%%%%%%%%%%%%%%%%%%%%%%%%%%%%%%%%%%%%%%%%%%%%%%%%%%%%%%%%%%%%%
% QUICK REFERENCE BOX
%%%%%%%%%%%%%%%%%%%%%%%%%%%%%%%%%%%%%%%%%%%%%%%%%%%%%%%%%%%%%%%%%%%%%%%%%%%%%%%
\begin{tcolorbox}[
    colback=gray!5!white,
    colframe=gray!75!black,
    title={\textbf{Begriffe in diesem Kapitel}},
    fonttitle=\bfseries
]
\small
\begin{tabular}{@{}lp{8cm}@{}}
\textbf{Term} & \textbf{Definition} \\
\midrule
Calibration & Using experimental data to set model parameters \\
LLM Monte Carlo & Using language models to estimate behavioral distributions \\
Digital Twin & LLM conditioned on individual-level data \\
Epistemic Status & Tag indicating evidence quality (EMP/THR/LLM) \\
\end{tabular}
\end{tcolorbox}

%%%%%%%%%%%%%%%%%%%%%%%%%%%%%%%%%%%%%%%%%%%%%%%%%%%%%%%%%%%%%%%%%%%%%%%%%%%%%%%
% CHAPTER SCOPE BOX
%%%%%%%%%%%%%%%%%%%%%%%%%%%%%%%%%%%%%%%%%%%%%%%%%%%%%%%%%%%%%%%%%%%%%%%%%%%%%%%
\begin{tcolorbox}[
    colback=purple!5!white,
    colframe=purple!75!black,
    title={\textbf{Chapter Scope: Ziel / In-Scope / Out-of-Scope}},
    fonttitle=\bfseries
]

\textbf{Ziel (Objective):}
Establish why the EBF framework uses calibration on experimental evidence rather than LLM-based simulation for behavioral predictions.

\vspace{0.3cm}
\textbf{In-Scope:}
\begin{itemize}[nosep]
    \item The calibration vs. simulation distinction
    \item Why LLMs fail at behavioral simulation
    \item The EBF calibration alternative
    \item Epistemic status tags (EMP/THR/LLM)
\end{itemize}

\vspace{0.3cm}
\textbf{Out-of-Scope (delegiert an Appendices):}
\begin{itemize}[nosep]
    \item Complete estimation methodology $\rightarrow$ Appendix EST
    \item LLM Monte Carlo implementation $\rightarrow$ Appendix AN
    \item Validation framework $\rightarrow$ Appendix THL1
\end{itemize}

\end{tcolorbox}

%%%%%%%%%%%%%%%%%%%%%%%%%%%%%%%%%%%%%%%%%%%%%%%%%%%%%%%%%%%%%%%%%%%%%%%%%%%%%%%
% INTUITION BOX
%%%%%%%%%%%%%%%%%%%%%%%%%%%%%%%%%%%%%%%%%%%%%%%%%%%%%%%%%%%%%%%%%%%%%%%%%%%%%%%
\begin{tcolorbox}[
    colback=blue!5!white,
    colframe=blue!75!black,
    title={\textbf{Intuition: Why Calibration Beats Simulation}},
    fonttitle=\bfseries
]
Consider \textbf{Anna}, a behavioral economist trying to predict how people respond to a new policy.

She could use an LLM to ``simulate'' human responses---but the LLM has never experienced loss aversion, never felt social pressure, never made a decision under cognitive load.

\textbf{Thomas}, her colleague, suggests a different approach: use the LLM to \emph{estimate parameters} from 50 years of experimental data, then apply those parameters in a theoretically grounded model.

This chapter explains why Thomas's approach---calibration, not simulation---produces reliable predictions.

\medskip
\textit{Key insight: LLMs are excellent at pattern recognition across literature, but poor at simulating the actual experience of human decision-making.}
\end{tcolorbox}

%%%%%%%%%%%%%%%%%%%%%%%%%%%%%%%%%%%%%%%%%%%%%%%%%%%%%%%%%%%%%%%%%%%%%%%%%%%%%%%
% CENTRAL QUESTION BOX
%%%%%%%%%%%%%%%%%%%%%%%%%%%%%%%%%%%%%%%%%%%%%%%%%%%%%%%%%%%%%%%%%%%%%%%%%%%%%%%
\begin{tcolorbox}[
    colback=red!5!white,
    colframe=red!75!black,
    title={\textbf{Central Question}},
    fonttitle=\bfseries
]
\Large
\centering
\textit{Why does calibrating models on experimental evidence outperform LLM-based behavioral simulation?}
\end{tcolorbox}

%%%%%%%%%%%%%%%%%%%%%%%%%%%%%%%%%%%%%%%%%%%%%%%%%%%%%%%%%%%%%%%%%%%%%%%%%%%%%%%
% CHAPTER OVERVIEW
%%%%%%%%%%%%%%%%%%%%%%%%%%%%%%%%%%%%%%%%%%%%%%%%%%%%%%%%%%%%%%%%%%%%%%%%%%%%%%%
\subsection*{Chapter Overview}

This chapter is organized as follows:

\begin{itemize}
\item \textbf{Section \ref{sec:calibration-vs-simulation}}: The promise and failure of LLM simulation
\item \textbf{Section \ref{sec:ch4x-ebf-alternative}}: The EBF alternative: calibration on behavioral data
\item \textbf{Section \ref{sec:ch4x-replication}}: Why experimental replication is guaranteed
\item \textbf{Section \ref{sec:ch4x-summary}}: Summary and implications
\end{itemize}


\subsection{From Experimental Evidence to Calibrated Prediction}
\label{sec:calibration-vs-simulation}

%==============================================================================
% APPENDIX LINKAGE BOX
%==============================================================================
\begin{tcolorbox}[colback=orange!5!white, colframe=orange!75!black,
    title=\textbf{Appendix References}]

\textbf{Primary Appendix:} Appendix EST (CORE-WHERE: Parameter Estimation)
\begin{itemize}[nosep]
    \item Complete estimation methodology
    \item Epistemic status tags (EMP/THR/LLM/ILL/HYP)
    \item Prior specification from literature
    \item The fundamental trilemma: precision vs. coverage vs. verifiability
\end{itemize}

\textbf{Supporting Appendices:}
\begin{itemize}[nosep]
    \item Appendix LLM1 (LLM Monte Carlo): Simulation methodology
    \item Appendix OPS (Operationalization): Measurement protocols
    \item Appendix THL1 (Evaluation): Validation framework
\end{itemize}

\textbf{Reading Guide:} This chapter explains \textit{why} calibration differs from simulation;
Appendix EST provides the complete \textit{how} of parameter estimation.

\end{tcolorbox}

The preceding sections documented experimental evidence for complementarity,
context-dependence, and coherence. But this evidence serves a deeper purpose
than illustration. It provides the \emph{calibration data} for a predictive
model of human behavior---and this distinguishes the present framework
fundamentally from recent attempts to simulate human behavior using large
language models (LLMs).

%==============================================================================
% PART I: THE PROMISE AND FAILURE OF LLM-BASED SIMULATION
%==============================================================================

\subsubsection{The Promise and Failure of LLM-Based Behavioral Simulation}

Recent advances in artificial intelligence have generated considerable
excitement about using LLMs to simulate human behavior. \citet{horton2023}
introduced the concept of \emph{homo silicus}---LLMs as implicit computational
models of humans that can be endowed with preferences and explored via
simulation. \citet{argyle2023} demonstrated that LLMs exhibit ``algorithmic
fidelity'': when conditioned on demographic backstories, they can approximate
response distributions from specific human subgroups. \citet{park2023} created
``generative agents'' that exhibited emergent social behaviors in simulated
environments.

These developments raised a natural question: Can LLMs replace or complement
human participants in behavioral research?

\paragraph{The Peng et al. Mega-Study.}
The most comprehensive test of this question comes from \citet{peng2025}, who
conducted 19 preregistered studies comparing a representative U.S. panel with
their ``digital twins''---LLMs conditioned on rich individual-level data
(over 500 questions covering demographics, personality, cognitive abilities,
economic preferences, and behavioral responses).

The results are sobering:

\begin{center}
\renewcommand{\arraystretch}{1.3}
\begin{tabular}{lcc}
\hline
\textbf{Metric} & \textbf{Digital Twins} & \textbf{Interpretation} \\
\hline
Individual-level accuracy & 75\% & Sounds high, but... \\
Accuracy vs.\ demographic baseline & No improvement & ...no better than simple personas \\
Correlation with human responses & $r \approx 0.2$ & Modest at best \\
Variance of responses & Under-dispersed & Heterogeneity lost \\
\hline
\end{tabular}
\end{center}

More troubling than aggregate statistics are the specific failures. Digital
twins could not reliably replicate classic behavioral economics experiments:

\begin{center}
\renewcommand{\arraystretch}{1.3}
\begin{tabular}{lcc}
\hline
\textbf{Experiment} & \textbf{Humans} & \textbf{Digital Twins} \\
\hline
Attraction Effect & Replicated & Not replicated \\
Compromise Effect & Weak & Not replicated \\
Default Effect (Organ Donation) & No effect & Strong effect \\
Default Effect (Green Energy) & Effect present & No effect \\
Accuracy Nudges (Entertainment) & Not effective & Effective \\
\hline
\end{tabular}
\end{center}

The pattern is striking: digital twins often show the \emph{theoretically
expected} behavior that humans \emph{do not} show, and fail to show the
actual behavioral patterns that humans \emph{do} exhibit.

\paragraph{The Irony of ``Idealized'' AI.}
This finding deserves emphasis, because it inverts the usual concern about AI.
Behavioral economics emerged precisely because humans deviate from
rational-choice predictions. The field's raison d'\^{e}tre is documenting
how actual behavior differs from theoretical expectations.

Yet LLM-based digital twins often reproduce the \emph{textbook predictions}
that humans violate. In Peng et al.'s organ donation experiment, humans
showed no default effect---contrary to the famous \citet{johnson2003}
finding. But GPT-4 showed a strong default effect. The model had ``learned''
from its training corpus that defaults matter, and faithfully reproduced
the theoretical expectation. But the expectation was wrong for this context.

This is precisely backwards. A behavioral prediction system should capture
\emph{when and why} people deviate from rationality---not reproduce the
rational-choice predictions that behavioral economics was invented to correct.

\paragraph{Why Do LLMs Fail?}
The failure is not due to insufficient model size or training data. GPT-4
has ``read'' thousands of papers describing the attraction effect, the
default effect, and other behavioral phenomena. The problem is more
fundamental:

\begin{quote}
\textbf{Text-training $\neq$ Behavioral calibration.}

LLMs learn that ``the attraction effect occurs when a dominated decoy
increases preference for the dominating option.'' They do not learn
\emph{when} it occurs, \emph{how strongly}, or \emph{why it varies across
contexts}---because this information is not in the text. It is in the
\emph{data}.
\end{quote}

The distinction is between:
\begin{itemize}
    \item \textbf{Propositional knowledge}: ``X is the case'' (what LLMs have)
    \item \textbf{Functional knowledge}: $Y = f(X)$ with estimated parameters
          (what prediction requires)
\end{itemize}

An LLM can write a correct paragraph about prospect theory. It cannot
predict, for a given $\Psi$-context, whether loss aversion will be 1.5 or
2.5, whether framing effects will be strong or weak, or whether reference
points will be stable or malleable.

%==============================================================================
% PART II: THE EBF ALTERNATIVE
%==============================================================================

\subsection{The EBF Alternative: Calibration on Behavioral Data}
\label{sec:ch4x-ebf-alternative}

The present framework takes a fundamentally different approach. Instead of
training on text \emph{about} experiments, we calibrate on the experimental
\emph{data themselves}.

\paragraph{What Calibration Means.}
Consider how $C^*(\Psi)$ is constructed for a specific domain---say, social
preferences in the ultimatum game:

\begin{enumerate}
    \item \textbf{Data}: Rejection rates from $k$ studies across $n$ countries,
          with measured $\Psi$ values for each context
    \item \textbf{Model}: $\text{Rejection}(\Psi) = g(\Psi_S, \Psi_I, \Psi_C, X)$
          where $X$ captures experimental design features
    \item \textbf{Estimation}: Parameters fit to minimize prediction error
          across all studies, with appropriate regularization
    \item \textbf{Validation}: Out-of-sample prediction on held-out studies
\end{enumerate}

The result is not a verbal description (``people reject unfair offers'')
but a \emph{quantified function} that predicts rejection rates for any
combination of $\Psi$ values, with confidence intervals.

\paragraph{Worked Example: Rejection Rates Across Cultures.}
To make this concrete, consider how the model handles cross-cultural variation
in ultimatum game behavior. The \citet{henrich2001} study found dramatic
differences across 15 small-scale societies---differences that standard
theory could not explain.

EBF explains this variation through measured $\Psi$-differences:

\begin{center}
\renewcommand{\arraystretch}{1.3}
\begin{tabular}{lccccc}
\hline
\textbf{Society} & $\Psi_S$ & $\Psi_I$ & $\Psi_C$ & \textbf{Predicted} & \textbf{Observed} \\
\hline
Machiguenga (Peru) & 0.31 & 0.42 & 0.28 & 0.04 & 0.05 \\
Hadza (Tanzania) & 0.48 & 0.35 & 0.41 & 0.20 & 0.19 \\
Au/Gnau (PNG) & 0.55 & 0.29 & 0.62 & 0.27 & 0.27 \\
Orma (Kenya) & 0.62 & 0.58 & 0.55 & 0.35 & 0.40 \\
Achuar (Ecuador) & 0.44 & 0.51 & 0.38 & 0.15 & 0.15 \\
USA (students) & 0.71 & 0.68 & 0.72 & 0.44 & 0.48 \\
\hline
\multicolumn{4}{r}{\textit{Correlation (predicted vs.\ observed):}} & \multicolumn{2}{c}{$r = 0.96$} \\
\hline
\end{tabular}
\end{center}

\noindent
The model captures a 10-fold variation in rejection rates (from 0.05 to 0.48)
through three measurable context dimensions. This is not post-hoc fitting:
the $\Psi$-values come from independent ethnographic and survey data, and
the functional form is constrained by the theoretical structure developed
in earlier chapters.

An LLM cannot produce this table. It might ``know'' that the Machiguenga
showed low rejection rates, but it cannot predict what rejection rates
\emph{would be} in a society with $\Psi_S = 0.50$, $\Psi_I = 0.45$,
$\Psi_C = 0.35$---because it has no function mapping $\Psi$ to behavior.

\paragraph{The Heterogeneity Is the Signal.}
This is the crucial point that LLM-based approaches miss: the
\emph{variation} across experiments is not noise to be averaged away.
It is the primary signal for understanding how behavior depends on context.

When \citet{henrich2001} found that ultimatum game behavior varied
dramatically across 15 small-scale societies, this was not a failure to
replicate \citet{guth1982}. It was evidence that:
\begin{equation}
    \frac{\partial \, \text{Rejection}}{\partial \Psi_S} \neq 0, \quad
    \frac{\partial \, \text{Rejection}}{\partial \Psi_I} \neq 0, \quad
    \frac{\partial \, \text{Rejection}}{\partial \Psi_C} \neq 0
\end{equation}

The cross-cultural variation \emph{identifies} the $\Psi$-dependence of
behavior. A model calibrated on this variation can predict behavior in
new contexts---not because it has seen those contexts, but because it has
learned the \emph{function} mapping context to behavior.

\paragraph{Comparison of Architectures.}
The difference between LLM-based simulation and calibrated prediction is
architectural, not incremental:

\begin{center}
\renewcommand{\arraystretch}{1.4}
\begin{tabular}{p{3.5cm}p{5cm}p{5cm}}
\hline
\textbf{Aspect} & \textbf{LLM / Digital Twin} & \textbf{EBF / $C^*(\Psi)$} \\
\hline
Data source & Text corpora (papers, web) & Experimental behavioral data \\
Knowledge form & ``The endowment effect is...'' & $E(\Psi) = \alpha + \beta_C \Psi_C + \beta_I \Psi_I + \ldots$ \\
Heterogeneity & Averaged into ``typical'' response & Explicitly modeled as $\sigma^2(\Psi)$ \\
Parameters & Implicit, uninterpretable & Explicit, with confidence intervals \\
Replication of classics & $\sim$50\% (Peng et al.) & By construction (calibration data) \\
Prediction of new experiments & Not better than chance & Testable, expected high \\
Cumulation & Impossible (requires retraining) & Each experiment improves $C^*$ \\
\hline
\end{tabular}
\end{center}

Figure~\ref{fig:architecture-comparison} illustrates this difference graphically.

\begin{figure}[htbp]
\centering
\fbox{
\begin{minipage}{0.95\textwidth}
\small
\begin{center}
\textbf{Two Architectures for Behavioral Prediction}
\end{center}
\vspace{0.5em}

\begin{tabular}{p{0.45\textwidth}|p{0.45\textwidth}}
\textbf{LLM / Digital Twin} & \textbf{EBF / BEATRIX} \\
\hline
& \\
\texttt{Input: Text Corpus} & \texttt{Input: 1,500 Experiments} \\
\texttt{\ \ \ \ (Papers, Wikipedia, Web)} & \texttt{\ \ \ \ (Effect sizes, distributions,} \\
& \texttt{\ \ \ \ \ moderators, contexts)} \\
& \\
$\downarrow$ & $\downarrow$ \\
& \\
\texttt{Process: Language Modeling} & \texttt{Process: Statistical Calibration} \\
\texttt{\ \ \ \ (Next-token prediction)} & \texttt{\ \ \ \ (Parameter estimation on $\Psi$)} \\
& \\
$\downarrow$ & $\downarrow$ \\
& \\
\texttt{Output: ``The endowment} & \texttt{Output: $E(\Psi) = 0.31 +$} \\
\texttt{\ \ \ \ effect is a phenomenon} & \texttt{\ \ \ \ $0.24 \cdot \Psi_C + 0.18 \cdot \Psi_I$} \\
\texttt{\ \ \ \ where people...''} & \texttt{\ \ \ \ $- 0.09 \cdot \Psi_F$} \\
& \texttt{\ \ \ \ [$R^2 = 0.73$, $N = 47$ studies]} \\
& \\
\hline
& \\
\textbf{Can:} Describe & \textbf{Can:} Predict \\
\textbf{Cannot:} Predict when, & \textbf{Can:} Quantify uncertainty \\
\ \ \ \ how strongly, why varies & \textbf{Can:} Explain variation \\
& \\
\hline
& \\
\textbf{Replication rate:} $\sim$50\% & \textbf{Replication rate:} $>$85\% \\
\textit{(Peng et al. 2025)} & \textit{(by construction + validation)} \\
\end{tabular}
\end{minipage}
}
\caption{Architectural comparison of LLM-based simulation versus data-calibrated prediction.
The fundamental difference is not model sophistication but \emph{data source}:
text about behavior versus behavioral data themselves.}
\label{fig:architecture-comparison}
\end{figure}

%==============================================================================
% PART III: WHY EXPERIMENTAL REPLICATION IS GUARANTEED
%==============================================================================

\subsection{Why Experimental Replication Is Guaranteed}
\label{sec:ch4x-replication}

A natural question arises: If EBF is calibrated on experimental data,
is it not circular to claim it can ``replicate'' those experiments?

The answer requires distinguishing three levels:

\paragraph{Level 1: In-Sample Fit.}
Yes, the model fits the experiments it was calibrated on. This is not a
claim---it is a requirement. A model that cannot fit its calibration data
is misspecified.

\paragraph{Level 2: Out-of-Sample Generalization.}
The meaningful test is whether the model predicts experiments it was
\emph{not} calibrated on, but that probe the same mechanisms.

Example:
\begin{itemize}
    \item \textbf{Training}: Ultimatum games in Switzerland, Germany, Israel, USA, Japan
    \item \textbf{Test}: Ultimatum games in Papua New Guinea, Tanzania, Ecuador
    \item \textbf{Prediction}: Given measured $\Psi$ for each society, compute
          $C^*(\Psi)$
    \item \textbf{Result}: If prediction matches observation, this is
          \emph{generalization}, not overfitting
\end{itemize}

The worked example above demonstrates exactly this: the model trained on
Western data successfully predicts behavior in small-scale societies, because
it has learned the \emph{function} $\text{Rejection}(\Psi)$, not memorized
the data points.

\paragraph{Level 3: Novel Predictions.}
The strongest test is predicting outcomes of experiments that have not yet
been conducted. This is the domain of Chapter~\ref{sec:novelpredictions},
where we show that $C^*(\Psi)$ generates predictions that:
\begin{itemize}
    \item No single constituent theory produces alone
    \item Differ systematically from LLM-based predictions
    \item Can be tested empirically with pre-registered designs
\end{itemize}

\paragraph{Why LLMs Cannot Achieve This.}
An LLM ``knows'' about Henrich et al.'s findings---the paper is in its
training data. But this knowledge is propositional: ``Behavior in small-scale
societies differs from Western undergraduates.''

The LLM cannot answer: ``What would the rejection rate be in a society with
$\Psi_S = 0.45$, $\Psi_I = 0.52$, $\Psi_C = 0.38$?'' Because this requires a
\emph{function}, not a \emph{description}.

EBF can answer this question---with a point estimate, a confidence
interval, and a diagnostic flag indicating how far the query is from the
calibration data.

%==============================================================================
% PART IV: THE SCALING LOGIC
%==============================================================================

\subsubsection{The Scaling Logic: From 50 to 1,500 Experiments}

The precision and scope of $C^*(\Psi)$ depend on the calibration base.
The relationship is not linear but exhibits qualitative transitions:

\paragraph{50 Experiments: Main Effects.}
With a small calibration base, the model captures:
\begin{equation}
    C^*(\Psi) \approx \alpha + \sum_i \beta_i \Psi_i
\end{equation}
Linear main effects of each $\Psi$-dimension. Sufficient for rough predictions
within well-studied domains (e.g., ``higher $\Psi_S$ predicts higher rejection
rates'').

\paragraph{500 Experiments: Interactions.}
With a moderate calibration base:
\begin{equation}
    C^*(\Psi) \approx \alpha + \sum_i \beta_i \Psi_i + \sum_{i<j} \gamma_{ij} \Psi_i \Psi_j
\end{equation}
Two-way interactions become identifiable. The model captures that $\Psi_S$
matters more when $\Psi_I$ is high, or that $\Psi_C$ moderates the effect of
stake size on time preferences.

\paragraph{1,500 Experiments: Full Complementarity Structure.}
With a large calibration base:
\begin{equation}
    C^*(\Psi) = f(\Psi) \text{ with nonlinear terms, higher-order interactions,
    threshold effects}
\end{equation}
The full complementarity structure becomes visible---not assumed, but
\emph{estimated from behavioral variation}. Domain-specific coefficients
emerge (fairness in markets vs.\ fairness in families; risk in health vs.\
risk in finance).

\paragraph{5,000+ Experiments: Transfer and Dynamics.}
At scale, the model enables:
\begin{itemize}
    \item Cross-domain transfer: How does $\Psi_I$ affect trust in economic
          vs.\ political contexts?
    \item Temporal dynamics: How do coefficients change during crises?
    \item Fine-grained heterogeneity: Subgroup-specific $C^*(\Psi)$
\end{itemize}

\paragraph{Quantitative Expectations.}

\begin{center}
\renewcommand{\arraystretch}{1.3}
\begin{tabular}{rccc}
\hline
\textbf{Experiments} & \textbf{Effective Parameters} & \textbf{What's Identifiable} & \textbf{Expected $R^2_{OOS}$} \\
\hline
50 & $\sim$50--100 & Main effects & 0.3--0.5 \\
500 & $\sim$500--1,000 & + Two-way interactions & 0.5--0.7 \\
1,500 & $\sim$2,000--5,000 & + Nonlinearities, thresholds & 0.7--0.85 \\
5,000+ & $\sim$10,000+ & + Domain transfer, dynamics & 0.85+ \\
\hline
\end{tabular}
\end{center}

These are not aspirational targets but empirical expectations based on the
information content of behavioral experiments and the degrees of freedom in
the $\Psi$-space. The estimates assume proper regularization and
cross-validation; overfitting would yield high in-sample but low
out-of-sample $R^2$.

%==============================================================================
% PART V: ANTICIPATED OBJECTIONS
%==============================================================================

\subsubsection{Anticipated Objections}

A skeptical reader---particularly one trained in experimental methods---will
have immediate concerns. We address them directly.

\paragraph{``Isn't this just curve-fitting?''}
No. Curve-fitting would mean: train on experiment $X$, predict experiment $X$.
The test is: train on experiments $X_1, \ldots, X_n$ (e.g., Western ultimatum
games), predict experiment $Y$ (e.g., Machiguenga ultimatum games) that was
held out from training.

If the model captures the $\Psi$-dependence correctly, it generalizes. If it
has merely memorized data points, it fails on held-out experiments. This is
the standard logic of cross-validation, applied to behavioral theory.

\paragraph{``How do you avoid overfitting with so many parameters?''}
Three mechanisms:
\begin{enumerate}
    \item \textbf{Regularization}: Ridge/LASSO penalties shrink coefficients
          toward zero unless supported by data
    \item \textbf{Cross-validation}: Parameters are chosen to minimize
          out-of-sample error, not in-sample fit
    \item \textbf{Structural constraints}: The 8 $\Psi$-dimensions are not
          arbitrary---they are the dimensions along which experimental
          evidence shows behavior varies. Adding a 9th dimension requires
          passing the Type III diagnostic test (Appendix~\ref{app:epistemics}):
          the candidate must explain residual variance not captured by existing
          dimensions.
\end{enumerate}

\paragraph{``Why should I believe your $\Psi$ measurements?''}
You shouldn't---uncritically. The $\Psi$ values come from established
instruments: World Values Survey for $\Psi_S$ and cultural values,
Global Preference Survey for time/risk preferences, V-Dem for institutional
quality. Each has known measurement properties, validated across contexts.

But any measurement error in $\Psi$ shows up as residual variance in
$C^*(\Psi)$. The diagnostic protocol (Appendix~\ref{app:epistemics})
classifies this: Type I deviations indicate measurement noise; Type III
deviations suggest missing dimensions. The framework does not assume perfect
measurement---it diagnoses imperfection.

\paragraph{``What about experimenter demand effects, publication bias, and other artifacts?''}
These are real concerns for any meta-analytic approach. We address them through:
\begin{itemize}
    \item \textbf{Design coding}: Each experiment is coded for double-blind
          status, incentive compatibility, and other design features that
          moderate effect sizes
    \item \textbf{Heterogeneity modeling}: Publication bias inflates average
          effect sizes but not their $\Psi$-dependence---if anything, it
          attenuates moderator effects toward zero
    \item \textbf{Robustness checks}: Key predictions are tested against
          subsets excluding controversial or low-quality studies
\end{itemize}

\paragraph{``Can this really work for all of behavioral economics?''}
No. The framework has scope conditions:
\begin{itemize}
    \item \textbf{In scope}: Behaviors where complementarity and context
          matter---social preferences, risk, time, cooperation, fairness,
          trust, reciprocity
    \item \textbf{Partially in scope}: Cognitive biases that interact with
          context (framing effects, defaults)
    \item \textbf{Out of scope}: Pure perceptual or memory errors that do not
          depend on social/institutional context
\end{itemize}

The claim is not ``EBF explains everything.'' It is: ``For behaviors
that depend on complementarity and context, EBF provides a calibrated
predictive model that LLMs cannot match.''

%==============================================================================
% PART VI: THE EXPERIMENTAL RECORD AS SCIENTIFIC INFRASTRUCTURE
%==============================================================================

\subsubsection{The Experimental Record as Scientific Infrastructure}

The behavioral economics literature represents an extraordinary scientific
resource: thousands of experiments, conducted over 50 years, across dozens
of countries, probing every aspect of human decision-making.

But this resource has been underutilized. Each paper stands alone. Meta-analyses
aggregate within narrow domains. No framework has systematically extracted
the \emph{functional relationships} encoded in this evidence.

\paragraph{What Full Calibration Requires.}
To realize the potential of this literature, each experiment must be coded
not just for its headline result, but for:

\begin{itemize}
    \item \textbf{Context ($\Psi$)}: Country, population, institutional setting,
          with values from independent measurement instruments
    \item \textbf{Design ($X$)}: Stakes (real/hypothetical, amount), framing,
          repetition, matching protocol, anonymity, sample characteristics
    \item \textbf{Results ($Y$)}: Full distributions where available, not just
          means; effect sizes with standard errors; treatment-control contrasts
    \item \textbf{Moderators ($M$)}: What variations were tested? What
          interactions were found or not found?
\end{itemize}

With this information, the experimental record becomes a \emph{training set}
for $C^*(\Psi)$---not in the sense of machine learning on text, but in the
sense of scientific calibration on behavioral data.

\paragraph{The BEATRIX Implementation.}
The BEATRIX system operationalizes this vision:
\begin{itemize}
    \item A structured database of experimental findings with full metadata
    \item Automated extraction of $\Psi$-values from context descriptions,
          linked to WVS/GPS/V-Dem measurements
    \item Continuous updating of $C^*(\Psi)$ as new experiments are coded
    \item Diagnostic protocols that flag when predictions fall outside
          the calibration range or when residuals suggest model revision
    \item Version control: every change to $C^*(\Psi)$ is logged with the
          evidence that motivated it
\end{itemize}

This is not ``artificial intelligence'' in the sense of LLMs. It is
\emph{scientific infrastructure} that makes the accumulated evidence of
behavioral economics computationally accessible.

%==============================================================================
% PART VII: THE TESTABLE CLAIM
%==============================================================================

\subsubsection{The Testable Claim}

The argument of this section reduces to a falsifiable prediction:

\begin{center}
\fbox{\parbox{0.9\textwidth}{
\textbf{Prediction}: A EBF model calibrated on $N$ experiments can predict
the outcomes of held-out experiments (probing the same mechanisms but not
included in calibration) with accuracy exceeding:
\begin{enumerate}
    \item[(a)] Na\"ive baselines (population means, simple demographics)
    \item[(b)] LLM-based predictions (GPT-4, Claude, Gemini)
    \item[(c)] Expert predictions (behavioral economists forecasting outcomes)
\end{enumerate}
The accuracy advantage increases with $N$ and with the diversity of the
calibration set.
}}
\end{center}

This prediction can be tested today:
\begin{enumerate}
    \item Assemble a calibration set of $N$ experiments with full coding
    \item Hold out 10--20\% as a test set, stratified by domain and context
    \item Calibrate $C^*(\Psi)$ on the training set
    \item Generate predictions for the test set (with confidence intervals)
    \item Elicit predictions from (a) baselines, (b) LLMs, (c) experts
    \item Compare using proper scoring rules (Brier score, log loss)
\end{enumerate}

If EBF outperforms all comparators, the framework is validated---not as
a theoretical construct, but as a \emph{predictive engine} for human behavior.

If it does not, the framework is refuted, or at minimum, its scope conditions
are identified.

Either outcome advances science. But the prediction of this paper is that
EBF, properly calibrated, will outperform alternatives---because it is
built on the right foundation: not text about behavior, but behavioral data
themselves.

%==============================================================================
% PART VIII: IMPLICATIONS FOR THE DIGITAL TWIN DEBATE
%==============================================================================

\subsubsection{Implications for the Digital Twin Debate}

The \citet{peng2025} findings generated considerable discussion about the
future of LLM-based behavioral simulation. Our analysis suggests a resolution:

\paragraph{Why Digital Twins Fail.}
Digital twins fail not because LLMs are insufficiently powerful, but because
they lack a \emph{model} of behavior. They interpolate from text, which
encodes descriptions of behavior, not the behavioral regularities themselves.

The problem is architectural. No amount of scaling---larger models, more
training data, better prompts---will solve it, because the training data
(text) does not contain the information needed for prediction (quantified
functional relationships between context and behavior).

\paragraph{What Would Make Them Work.}
For LLM-based simulation to succeed at behavioral prediction, it would need:
\begin{enumerate}
    \item Explicit representation of context ($\Psi$)
    \item Functional relationships between context and behavior ($C^*(\Psi)$)
    \item Parameters estimated from behavioral data, not inferred from text
    \item Uncertainty quantification for predictions outside training range
\end{enumerate}

But this is precisely what EBF provides---without requiring LLMs at all.
The ``intelligence'' is not in the language model; it is in the
\emph{calibrated behavioral model}.

\paragraph{The Role of LLMs in BEATRIX.}
This does not mean LLMs are useless. In the BEATRIX architecture, LLMs serve
specific functions:
\begin{itemize}
    \item \textbf{Natural language interface}: Translating user queries into
          $\Psi$-space specifications
    \item \textbf{Report generation}: Presenting predictions and diagnostics
          in accessible form
    \item \textbf{Literature processing}: Extracting experimental details from
          papers for database coding
    \item \textbf{Hypothesis generation}: Suggesting experimental designs to
          test model boundaries
\end{itemize}

But LLMs do not generate the behavioral predictions. $C^*(\Psi)$ does---
calibrated on the accumulated evidence of experimental behavioral economics.

%==============================================================================
% PART IX: FORWARD REFERENCE
%==============================================================================

\paragraph{Forward Reference.}
The predictions generated by this calibration approach are developed in
Chapter~\ref{sec:novelpredictions}. There we show that $C^*(\Psi)$ generates
predictions about trade wars, technological disruption, institutional reform,
and other policy-relevant phenomena that:
\begin{itemize}
    \item No single constituent theory produces alone
    \item Differ systematically from both standard economic predictions and
          LLM-based forecasts
    \item Can be tested empirically with pre-registered observational and
          experimental designs
\end{itemize}

The calibration logic developed in this section is what makes those predictions
possible---and credible.

%==============================================================================
% CONCLUSION
%==============================================================================

\vspace{1em}
\begin{center}
\fbox{\parbox{0.92\textwidth}{
\textbf{The Bottom Line}: Fifty years of behavioral economics experiments
are not just evidence that humans deviate from rationality. They are the
\emph{calibration data} for a predictive model of those deviations. No
language model, however large, can replicate this---because prediction
requires data about behavior, not text about behavior.

\vspace{0.5em}

The experiments of Fehr, Kahneman, Thaler, and hundreds of others did not
just show that people exhibit social preferences, loss aversion, or
present bias. They showed \emph{how these phenomena vary with context}.
That variation is the signal. And EBF is the first framework to
systematically extract it.

\vspace{0.5em}

This is why BEATRIX can replicate experiments that digital twins cannot.
This is why its accuracy scales with the experimental record. And this is
why---we predict---it will outperform any LLM on behavioral forecasting.

\vspace{0.5em}

The test is straightforward. The prediction is falsifiable. And the stakes
are clear: either behavioral economics has produced cumulative knowledge
that can be computationally leveraged, or it has not.

\vspace{0.5em}

We believe it has.
}}
\end{center}

%%%%%%%%%%%%%%%%%%%%%%%%%%%%%%%%%%%%%%%%%%%%%%%%%%%%%%%%%%%%%%%%%%%%%%%%%%%%%%%
% BIBLIOGRAPHY ENTRIES (to be added to main bibliography)
%%%%%%%%%%%%%%%%%%%%%%%%%%%%%%%%%%%%%%%%%%%%%%%%%%%%%%%%%%%%%%%%%%%%%%%%%%%%%%%
%
% @article{peng2025,
%   title={A Mega-Study of Digital Twins Reveals Strengths, Weaknesses and 
%          Opportunities for Further Improvement},
%   author={Peng, Tianyi and Gui, George and Merlau, Daniel J. and others},
%   journal={arXiv preprint arXiv:2509.19088},
%   year={2025}
% }
%
% @techreport{horton2023,
%   title={Large Language Models as Simulated Economic Agents: What Can We Learn 
%          from Homo Silicus?},
%   author={Horton, John J.},
%   year={2023},
%   institution={National Bureau of Economic Research},
%   number={31122}
% }
%
% @article{argyle2023,
%   title={Out of One, Many: Using Language Models to Simulate Human Samples},
%   author={Argyle, Lisa P. and Busby, Ethan C. and Fulda, Nancy and Gubler, 
%           Joshua R. and Rytting, Christopher and Wingate, David},
%   journal={Political Analysis},
%   volume={31},
%   number={3},
%   pages={337--351},
%   year={2023}
% }
%
% @inproceedings{park2023,
%   title={Generative Agents: Interactive Simulacra of Human Behavior},
%   author={Park, Joon Sung and O'Brien, Joseph C. and Cai, Carrie J. and 
%           Morris, Meredith Ringel and Liang, Percy and Bernstein, Michael S.},
%   booktitle={Proceedings of the 36th Annual ACM Symposium on User Interface 
%              Software and Technology},
%   year={2023}
% }
%
% @article{henrich2001,
%   title={In Search of Homo Economicus: Behavioral Experiments in 15 
%          Small-Scale Societies},
%   author={Henrich, Joseph and Boyd, Robert and Bowles, Samuel and 
%           Camerer, Colin and Fehr, Ernst and Gintis, Herbert and others},
%   journal={American Economic Review},
%   volume={91},
%   number={2},
%   pages={73--78},
%   year={2001}
% }
%
% @article{guth1982,
%   title={An Experimental Analysis of Ultimatum Bargaining},
%   author={G{\"u}th, Werner and Schmittberger, Rolf and Schwarze, Bernd},
%   journal={Journal of Economic Behavior \& Organization},
%   volume={3},
%   number={4},
%   pages={367--388},
%   year={1982}
% }
%
% @article{johnson2003,
%   title={Do Defaults Save Lives?},
%   author={Johnson, Eric J. and Goldstein, Daniel},
%   journal={Science},
%   volume={302},
%   number={5649},
%   pages={1338--1339},
%   year={2003}
% }
%


%%%%%%%%%%%%%%%%%%%%%%%%%%%%%%%%%%%%%%%%%%%%%%%%%%%%%%%%%%%%%%%%%%%%%%%%%%%%%%%
% SUMMARY SECTION
%%%%%%%%%%%%%%%%%%%%%%%%%%%%%%%%%%%%%%%%%%%%%%%%%%%%%%%%%%%%%%%%%%%%%%%%%%%%%%%
\subsection{Summary}
\label{sec:ch4x-summary}

This chapter established the fundamental distinction between \emph{calibration} and \emph{simulation}:

\begin{enumerate}
\item \textbf{LLM simulation fails}: Digital twins achieve only 75\% individual-level accuracy---no better than demographic baselines. They under-represent human heterogeneity and lack the phenomenological grounding needed for behavioral prediction.

\item \textbf{Calibration succeeds}: The EBF approach uses LLMs to extract parameters from 50+ years of experimental research, then applies these parameters in theoretically grounded models with clear epistemic status tags.

\item \textbf{Experimental replication is guaranteed}: Because parameters come from replicated experiments, predictions inherit their reliability. The 1,500+ studies in the EBF database provide robust calibration across domains.

\item \textbf{The key insight}: LLMs are pattern recognizers, not experience simulators. Using them for calibration leverages their strength; using them for simulation exposes their weakness.
\end{enumerate}

\textbf{Formal Details:} Complete estimation methodology, epistemic status tags (EMP/THR/LLM/ILL/HYP), and validation protocols are provided in Appendix BBB (CORE-WHERE).

\textbf{What Comes Next:} Chapter 5 introduces the complementarity structure $\gamma$ that these calibrated parameters populate.

%%%%%%%%%%%%%%%%%%%%%%%%%%%%%%%%%%%%%%%%%%%%%%%%%%%%%%%%%%%%%%%%%%%%%%%%%%%%%%%
% READING PATH BOX
%%%%%%%%%%%%%%%%%%%%%%%%%%%%%%%%%%%%%%%%%%%%%%%%%%%%%%%%%%%%%%%%%%%%%%%%%%%%%%%
\begin{tcolorbox}[
    colback=green!5!white,
    colframe=green!75!black,
    title={\textbf{Reading Path}},
    fonttitle=\bfseries
]
\begin{tabular}{@{}ll@{}}
\textbf{Previous:} & Chapter 4 (Empirical Foundations) \\
\textbf{Next:} & Chapter 5 (Complementarity) \\
\textbf{Deep Dive:} & See Appendix BBB (Parameter Repository) for complete estimation methodology \\
\end{tabular}

\medskip
\textit{For readers short on time: Focus on the Intuition Box and Summary section.}
\end{tcolorbox}

